\section{The \texttt{Lens\_elliptical} McXtrace Component}
X-ray compound refractive lens (CRL) with an elliptic profile

\subsection*{Identification}
\begin{itemize}
  \item \textbf{Author:} Jana Baltser and Erik Knudsen
  \item \textbf{Origin:} NBI
  \item \textbf{Date:} August 2010
\end{itemize}

\subsection*{Description}
A simple X-ray compound refractive lens (CRL) with an elliptic profile simulates the photons' movement on passing through it. Attenuation coefficient mu is taken from the NIST database and Be.txt

\subsection*{Input parameters}
Parameters in \textbf{boldface} are required; the others are optional.

\begin{longtable}{p{0.22\textwidth}p{0.12\textwidth}p{0.46\textwidth}p{0.14\textwidth}}
\toprule
\textbf{Name} & \textbf{Unit} & \textbf{Description} & \textbf{Default} \\
\midrule
\endhead
material\_datafile & Be.txt & File where the material parameters for the filter may be found. Format is similar to what may be found off the NIST website. & "Be.txt" \\
r1 & m & Radius of the profile along the X axis & 0.42e-3 \\
r2 & m & Radius of the profile along the Y axis & 0.8e-3 \\
w & m & Parabola parameter, constraining it along the propagation axis & 0.46e-3 \\
d & m & Distance between two surfaces of the lens along the propagation axis & 0.2e-4 \\
Transmission &  &  & 1 \\
N & 1 & Amount of single lenses in a stack & 1 \\
\bottomrule
\end{longtable}

\subsection*{Links}
\begin{itemize}
  \item Component source code found in file \texttt{Lens\_elliptical.comp}.
  \item material datafile obtained from http://physics.nist.gov/cgi-bin/ffast/ffast.pl
\end{itemize}
\IfFileExists{optics/Lens_elliptical_static.tex}{\input{optics/Lens_elliptical_static.tex}}{}